% assets/w05-gnc-loops.svg — 5주차 GNC 두 고리 (바깥 = 유도, 안쪽 = 제어)
%
%   bash _tools/tikz2svg.sh assets/src/w05-gnc-loops.tex assets/w05-gnc-loops.svg
%   bash _tools/svg2png.sh  assets/w05-gnc-loops.svg
%
% 이 그림이 말하는 것 다섯 가지
%   1) 바깥 고리(유도)가 내는 것은 **힘이 아니라 지령**이다 — psi_ref 와 u_ref
%   2) 바깥 고리의 출력이 그대로 안쪽 고리의 입력이다 (그래서 두 고리가 맞물린다)
%   3) 안쪽 고리는 3·4주차에 이미 만들었다. 이번 주차에 새로 만드는 칸은 왼쪽 하나다
%   4) 배분이 내는 것은 좌우 추력 F_L, F_R 이고, 프로펠러 법칙은 그 **다음** 칸이다
%   5) 되먹임이 둘이다 — 안쪽은 빠르고(u, psi, r) 바깥은 느리다(x_n, y_n, psi)
%
% 출처와 각색 (2026-10-14)
%   원본은 센서신호처리(대학원) 과목의 Otter 용 그림이다. 이 자료에 맞추어
%   고친 곳:
%     - 배가 Otter -> **WAM-V** (VRX). 추진기가 후방 2개, 좌우 반폭 half_beam
%     - 웨이포인트를 (N_k, E_k) -> **(x_k, y_k)**. 이 볼트에서 N 은 **요 모멘트**
%       전용이라 위치에 N 을 쓰면 한 기호가 두 가지 뜻이 된다 (표기 규칙)
%     - psi_d -> **psi_ref**, u_d -> **u_ref**. 모델의 신호 이름과 같게 맞춘다
%     - 적분 LOS 항(kappa * y_int)을 **뺀다**. 5주차 모델에 없는 항이다
%     - 배분 칸이 n 까지 내던 것을 갈라, 프로펠러 법칙을 **추진기** 칸으로 옮긴다.
%       모델의 Alloc / Thrusters 서브시스템 경계와 같아진다
%
% 좌표 메모 — 블록 중심과 폭 (mm)
%   gu (42,0) w44   hd (101,0) w52   su (97,-34) w46
%   al (154,-17) w34   th (198,-17) w32   sh (240,-17) w26
%   안쪽 되먹임 버스 y=-50, 바깥 되먹임 버스 y=-60
%   되먹임 수직선은 **블록 폭 밖**으로 지나게 잡았다 (x=124 는 su 오른끝 120 밖,
%   동시에 hd 폭 75..127 안이다). 이 값을 고치면 선이 상자를 관통한다

\documentclass[border=6pt]{standalone}
% capstone-style.tex — 이 볼트의 TikZ 그림이 공유하는 스타일
%
% 쓰는 법:  % capstone-style.tex — 이 볼트의 TikZ 그림이 공유하는 스타일
%
% 쓰는 법:  % capstone-style.tex — 이 볼트의 TikZ 그림이 공유하는 스타일
%
% 쓰는 법:  \input{capstone-style}   (assets/src/ 안에서)
% 빌드:     bash _tools/tikz2svg.sh assets/src/w02-node-topic.tex assets/w02-node-topic.svg
%
% 왜 TikZ 인가
%   손으로 SVG 좌표를 쓰면 화살촉·선 굵기·글자 크기가 그림마다 미묘하게 달라진다.
%   그것이 "자료가 어설퍼 보인다"의 정체다. 스타일을 한 번만 정하고 상대좌표로
%   배치하면 좌표를 고쳐도 그림이 무너지지 않는다.
%   수식은 본문과 같은 TeX 글꼴로 나온다.
%
% 한글
%   ko.TeX(DVI 경로)를 쓴다. XeLaTeX 은 TikZ 그래픽을 PDF special 로 내보내는데
%   dvisvgm 이 그것을 받으려면 Ghostscript < 10.01 이 필요하다. 이 PC 에는
%   10.07 이 깔려 있어 그 경로가 막힌다. latex -> DVI -> dvisvgm 이 유일하게
%   동작하는 조합이다.

\usepackage[hangul]{kotex}
\usepackage{tikz}
\usepackage{amsmath}
\usepackage{xcolor}
\usetikzlibrary{arrows.meta, positioning, calc, fit, backgrounds, decorations.pathmorphing}

% ---- 색 --------------------------------------------------------------------
% 강조는 하나만. 나머지는 검정. 흑백 인쇄에서도 읽혀야 한다.
\definecolor{ink}{HTML}{1A1A1A}
\definecolor{soft}{HTML}{666666}
\definecolor{accent}{HTML}{7C3AED}
\definecolor{warn}{HTML}{B91C1C}

% 계층 색 — Simulink 모델의 블록 색과 같은 규칙을 쓴다
\definecolor{guid}{HTML}{1D4ED8}    % 유도
\definecolor{ctrl}{HTML}{EA580C}    % 제어
\definecolor{thr}{HTML}{D97706}     % 추진기
\definecolor{plant}{HTML}{16A34A}   % 운동모델
\definecolor{nav}{HTML}{7C3AED}     % 항법 · 신호처리
\definecolor{pathc}{HTML}{0D9488}   % 계획 경로 · 항적
\definecolor{errc}{HTML}{C2410C}    % 오차
\definecolor{flow}{HTML}{0369A1}    % 조류 · 바람

% ---- 선과 화살촉 -----------------------------------------------------------
\tikzset{
  sig/.style   = {ink, line width=0.5pt, -{Stealth[length=4pt, width=3pt]}},
  wire/.style  = {ink, line width=0.5pt},
  fb/.style    = {ink, line width=0.5pt, densely dashed,
                  -{Stealth[length=4pt, width=3pt]}},
  asig/.style  = {accent, line width=0.5pt, densely dashed,
                  -{Stealth[length=4pt, width=3pt]}},
  % 블록 — 흰 바탕, 직각, 검은 테두리. 색은 테두리에만 준다
  blk/.style   = {draw=ink, line width=0.55pt, fill=white,
                  minimum width=18mm, minimum height=8.5mm, inner sep=2pt, font=\small},
  gblk/.style  = {blk, draw=guid},
  cblk/.style  = {blk, draw=ctrl},
  tblk/.style  = {blk, draw=thr},
  pblk/.style  = {blk, draw=plant},
  nblk/.style  = {blk, draw=nav},
  % 합산점 · 분기점
  sum/.style   = {draw=ink, line width=0.55pt, fill=white, circle,
                  inner sep=0pt, minimum size=4.6mm, font=\scriptsize},
  dot/.style   = {ink, fill=ink, circle, inner sep=0pt, minimum size=1.5mm},
  % 글자
  sname/.style = {font=\small, inner sep=1.5pt},
  note/.style  = {font=\scriptsize, soft, inner sep=1.5pt},
  wnote/.style = {font=\scriptsize, warn, inner sep=1.5pt},
  anote/.style = {font=\scriptsize, accent, inner sep=1.5pt},
  axis/.style  = {ink, line width=0.5pt},
  tick/.style  = {font=\scriptsize, soft},
  % 묶음 상자 — 서브시스템 표시
  grp/.style   = {draw=soft, line width=0.4pt, densely dotted, rounded corners=1mm,
                  inner sep=3mm},
}
   (assets/src/ 안에서)
% 빌드:     bash _tools/tikz2svg.sh assets/src/w02-node-topic.tex assets/w02-node-topic.svg
%
% 왜 TikZ 인가
%   손으로 SVG 좌표를 쓰면 화살촉·선 굵기·글자 크기가 그림마다 미묘하게 달라진다.
%   그것이 "자료가 어설퍼 보인다"의 정체다. 스타일을 한 번만 정하고 상대좌표로
%   배치하면 좌표를 고쳐도 그림이 무너지지 않는다.
%   수식은 본문과 같은 TeX 글꼴로 나온다.
%
% 한글
%   ko.TeX(DVI 경로)를 쓴다. XeLaTeX 은 TikZ 그래픽을 PDF special 로 내보내는데
%   dvisvgm 이 그것을 받으려면 Ghostscript < 10.01 이 필요하다. 이 PC 에는
%   10.07 이 깔려 있어 그 경로가 막힌다. latex -> DVI -> dvisvgm 이 유일하게
%   동작하는 조합이다.

\usepackage[hangul]{kotex}
\usepackage{tikz}
\usepackage{amsmath}
\usepackage{xcolor}
\usetikzlibrary{arrows.meta, positioning, calc, fit, backgrounds, decorations.pathmorphing}

% ---- 색 --------------------------------------------------------------------
% 강조는 하나만. 나머지는 검정. 흑백 인쇄에서도 읽혀야 한다.
\definecolor{ink}{HTML}{1A1A1A}
\definecolor{soft}{HTML}{666666}
\definecolor{accent}{HTML}{7C3AED}
\definecolor{warn}{HTML}{B91C1C}

% 계층 색 — Simulink 모델의 블록 색과 같은 규칙을 쓴다
\definecolor{guid}{HTML}{1D4ED8}    % 유도
\definecolor{ctrl}{HTML}{EA580C}    % 제어
\definecolor{thr}{HTML}{D97706}     % 추진기
\definecolor{plant}{HTML}{16A34A}   % 운동모델
\definecolor{nav}{HTML}{7C3AED}     % 항법 · 신호처리
\definecolor{pathc}{HTML}{0D9488}   % 계획 경로 · 항적
\definecolor{errc}{HTML}{C2410C}    % 오차
\definecolor{flow}{HTML}{0369A1}    % 조류 · 바람

% ---- 선과 화살촉 -----------------------------------------------------------
\tikzset{
  sig/.style   = {ink, line width=0.5pt, -{Stealth[length=4pt, width=3pt]}},
  wire/.style  = {ink, line width=0.5pt},
  fb/.style    = {ink, line width=0.5pt, densely dashed,
                  -{Stealth[length=4pt, width=3pt]}},
  asig/.style  = {accent, line width=0.5pt, densely dashed,
                  -{Stealth[length=4pt, width=3pt]}},
  % 블록 — 흰 바탕, 직각, 검은 테두리. 색은 테두리에만 준다
  blk/.style   = {draw=ink, line width=0.55pt, fill=white,
                  minimum width=18mm, minimum height=8.5mm, inner sep=2pt, font=\small},
  gblk/.style  = {blk, draw=guid},
  cblk/.style  = {blk, draw=ctrl},
  tblk/.style  = {blk, draw=thr},
  pblk/.style  = {blk, draw=plant},
  nblk/.style  = {blk, draw=nav},
  % 합산점 · 분기점
  sum/.style   = {draw=ink, line width=0.55pt, fill=white, circle,
                  inner sep=0pt, minimum size=4.6mm, font=\scriptsize},
  dot/.style   = {ink, fill=ink, circle, inner sep=0pt, minimum size=1.5mm},
  % 글자
  sname/.style = {font=\small, inner sep=1.5pt},
  note/.style  = {font=\scriptsize, soft, inner sep=1.5pt},
  wnote/.style = {font=\scriptsize, warn, inner sep=1.5pt},
  anote/.style = {font=\scriptsize, accent, inner sep=1.5pt},
  axis/.style  = {ink, line width=0.5pt},
  tick/.style  = {font=\scriptsize, soft},
  % 묶음 상자 — 서브시스템 표시
  grp/.style   = {draw=soft, line width=0.4pt, densely dotted, rounded corners=1mm,
                  inner sep=3mm},
}
   (assets/src/ 안에서)
% 빌드:     bash _tools/tikz2svg.sh assets/src/w02-node-topic.tex assets/w02-node-topic.svg
%
% 왜 TikZ 인가
%   손으로 SVG 좌표를 쓰면 화살촉·선 굵기·글자 크기가 그림마다 미묘하게 달라진다.
%   그것이 "자료가 어설퍼 보인다"의 정체다. 스타일을 한 번만 정하고 상대좌표로
%   배치하면 좌표를 고쳐도 그림이 무너지지 않는다.
%   수식은 본문과 같은 TeX 글꼴로 나온다.
%
% 한글
%   ko.TeX(DVI 경로)를 쓴다. XeLaTeX 은 TikZ 그래픽을 PDF special 로 내보내는데
%   dvisvgm 이 그것을 받으려면 Ghostscript < 10.01 이 필요하다. 이 PC 에는
%   10.07 이 깔려 있어 그 경로가 막힌다. latex -> DVI -> dvisvgm 이 유일하게
%   동작하는 조합이다.

\usepackage[hangul]{kotex}
\usepackage{tikz}
\usepackage{amsmath}
\usepackage{xcolor}
\usetikzlibrary{arrows.meta, positioning, calc, fit, backgrounds, decorations.pathmorphing}

% ---- 색 --------------------------------------------------------------------
% 강조는 하나만. 나머지는 검정. 흑백 인쇄에서도 읽혀야 한다.
\definecolor{ink}{HTML}{1A1A1A}
\definecolor{soft}{HTML}{666666}
\definecolor{accent}{HTML}{7C3AED}
\definecolor{warn}{HTML}{B91C1C}

% 계층 색 — Simulink 모델의 블록 색과 같은 규칙을 쓴다
\definecolor{guid}{HTML}{1D4ED8}    % 유도
\definecolor{ctrl}{HTML}{EA580C}    % 제어
\definecolor{thr}{HTML}{D97706}     % 추진기
\definecolor{plant}{HTML}{16A34A}   % 운동모델
\definecolor{nav}{HTML}{7C3AED}     % 항법 · 신호처리
\definecolor{pathc}{HTML}{0D9488}   % 계획 경로 · 항적
\definecolor{errc}{HTML}{C2410C}    % 오차
\definecolor{flow}{HTML}{0369A1}    % 조류 · 바람

% ---- 선과 화살촉 -----------------------------------------------------------
\tikzset{
  sig/.style   = {ink, line width=0.5pt, -{Stealth[length=4pt, width=3pt]}},
  wire/.style  = {ink, line width=0.5pt},
  fb/.style    = {ink, line width=0.5pt, densely dashed,
                  -{Stealth[length=4pt, width=3pt]}},
  asig/.style  = {accent, line width=0.5pt, densely dashed,
                  -{Stealth[length=4pt, width=3pt]}},
  % 블록 — 흰 바탕, 직각, 검은 테두리. 색은 테두리에만 준다
  blk/.style   = {draw=ink, line width=0.55pt, fill=white,
                  minimum width=18mm, minimum height=8.5mm, inner sep=2pt, font=\small},
  gblk/.style  = {blk, draw=guid},
  cblk/.style  = {blk, draw=ctrl},
  tblk/.style  = {blk, draw=thr},
  pblk/.style  = {blk, draw=plant},
  nblk/.style  = {blk, draw=nav},
  % 합산점 · 분기점
  sum/.style   = {draw=ink, line width=0.55pt, fill=white, circle,
                  inner sep=0pt, minimum size=4.6mm, font=\scriptsize},
  dot/.style   = {ink, fill=ink, circle, inner sep=0pt, minimum size=1.5mm},
  % 글자
  sname/.style = {font=\small, inner sep=1.5pt},
  note/.style  = {font=\scriptsize, soft, inner sep=1.5pt},
  wnote/.style = {font=\scriptsize, warn, inner sep=1.5pt},
  anote/.style = {font=\scriptsize, accent, inner sep=1.5pt},
  axis/.style  = {ink, line width=0.5pt},
  tick/.style  = {font=\scriptsize, soft},
  % 묶음 상자 — 서브시스템 표시
  grp/.style   = {draw=soft, line width=0.4pt, densely dotted, rounded corners=1mm,
                  inner sep=3mm},
}


\begin{document}
\begin{tikzpicture}[x=1mm, y=1mm]

% ===== 왼쪽 — 임무가 주는 것 (신호가 아니라 설정) ==========================
\node[blk, minimum width=24mm, align=center, font=\scriptsize] at (0,0)
      (wp) {\textbf{웨이포인트}\\$(x_k,\,y_k)$,\; $k=1\dots n$};
\node[blk, minimum width=24mm, align=center, font=\scriptsize] at (0,-34)
      (sp) {\textbf{목표 속도}\\$u_{\text{ref}} = 1.5$ m/s};

% ===== 바깥 고리 — 유도 ====================================================
\node[gblk, minimum width=44mm, minimum height=14mm, align=center, font=\scriptsize] at (42,0)
      (gu) {\textbf{유도}\; {\itshape outer loop}\\[2pt]
            $\psi_{\text{ref}} = \pi_p - \arctan\dfrac{y_e}{\Delta}$};

% ===== 안쪽 고리 — 제어 (3·4주차에 만든 것) ================================
\node[cblk, minimum width=52mm, minimum height=12mm, align=center, font=\scriptsize] at (101,0)
      (hd) {\textbf{헤딩 고리}\\[2pt]
            $N = K_{p\psi}\,\mathrm{ssa}(\psi_{\text{ref}}-\psi) - K_{d\psi}\,r$};
\node[cblk, minimum width=46mm, minimum height=12mm, align=center, font=\scriptsize] at (97,-34)
      (su) {\textbf{속도 고리}\\[2pt]
            $X = K_{pu}(u_{\text{ref}}-u) + K_{iu}\!\displaystyle\int\!(u_{\text{ref}}-u)\,dt$};

% ===== 배분 — X, N 을 좌우 추력으로 ========================================
\node[tblk, minimum width=34mm, minimum height=19mm, align=center, font=\scriptsize] at (154,-17)
      (al) {\textbf{배분}\\[2pt]
            $\begin{bmatrix} F_L \\ F_R \end{bmatrix}
             = \begin{bmatrix} \tfrac12 & \phantom{-}\tfrac{1}{2b} \\[2pt]
                               \tfrac12 & -\tfrac{1}{2b} \end{bmatrix}
               \begin{bmatrix} X \\ N \end{bmatrix}$\\[2pt]
            $b = \texttt{half\_beam} = 1.027$ m};

% ===== 추진기 — 프로펠러 법칙은 여기다 =====================================
\node[tblk, minimum width=32mm, minimum height=19mm, align=center, font=\scriptsize] at (198,-17)
      (th) {\textbf{추진기}\\[2pt]
            $n = \pm\sqrt{|F|/k}$\\[1pt]
            모터 지연 $\tau_n$\\[1pt]
            $F = k\,n\,|n|$};

% ===== 플랜트 ==============================================================
\node[pblk, minimum width=26mm, minimum height=19mm, align=center, font=\scriptsize] at (240,-17)
      (sh) {\textbf{WAM-V}\\운동 방정식\\[1pt]{\itshape 3자유도}};

% ===== 앞먹임 배선 ==========================================================
\draw[sig] (wp) -- (gu);
\draw[sig] (gu) -- node[sname, above, font=\scriptsize] {$\psi_{\text{ref}}$} (hd);
%  u_ref 이름표는 화살표 가운데가 아니라 끝쪽에 둔다. 가운데는 바깥 되먹임
%  수직선(x=42)이 지나는 자리라 글자와 선이 겹친다
\draw[sig] (sp) -- node[sname, above, pos=0.8, font=\scriptsize] {$u_{\text{ref}}$} (su);

\draw[sig] (hd.east) -| node[sname, pos=0.2, above, font=\scriptsize] {$N$} (al.north);
\draw[sig] (su.east) -| node[sname, pos=0.2, above, font=\scriptsize] {$X$} (al.south);
\draw[sig] (al) -- node[sname, above, font=\scriptsize] {$F_L, F_R$} (th);
\draw[sig] (th) -- node[sname, above, font=\scriptsize] {$F_L, F_R$} (sh);
\draw[sig] (sh.east) -- ++(9,0)
      node[sname, right, font=\scriptsize] {$x_n,\,y_n,\,\psi,\,u,\,v,\,r$};

% ===== 되먹임 ==============================================================
% 안쪽 — 빠르다. 제어기가 매 스텝 본다
\draw[wire] (sh.south) -- (240,-60);
\node[dot] at (240,-50) {};
\draw[wire] (240,-50) -- (88,-50);
\draw[sig]  (124,-50) -- (124,-6);
\draw[sig]  (88,-50)  -- (88,-40);
\node[sname, font=\scriptsize, anchor=west] at (125,-12) {$\psi,\;r$};
\node[sname, font=\scriptsize, anchor=east] at (86.5,-45) {$u$};
\node[wnote, anchor=south east, font=\scriptsize] at (238,-49.5) {안쪽 되먹임 — 빠름};

% 바깥 — 느리다. 위치가 웨이포인트에 닿았는지만 본다
\draw[wire] (240,-60) -- (42,-60);
\draw[sig]  (42,-60) -- (42,-7);
\node[note, above, font=\scriptsize] at (150,-60) {위치 $(x_n,\,y_n)$ 와 선수각 $\psi$};
\node[wnote, anchor=south east, font=\scriptsize] at (238,-59.5) {바깥 되먹임 — 느림};

% ===== 띠 제목 ==============================================================
\draw[wire, soft] (-14,15) -- (70,15);
\node[anchor=south west, font=\scriptsize\bfseries] at (-14,15.5)
      {바깥 고리\; \mdseries 유도는 \textbf{지령}을 낸다 — 힘을 직접 내지 않음};
\draw[wire, soft] (74,15) -- (253,15);
\node[anchor=south west, font=\scriptsize\bfseries, warn] at (74,15.5)
      {안쪽 고리\; \mdseries 3 · 4주차에 이미 만든 것};

% ===== 설명 =================================================================
\node[anote, anchor=south, font=\scriptsize] at (42,7.8) {이번 주차에 새로 만드는 칸};
\node[wnote, anchor=north, align=center, text width=26mm, font=\scriptsize] at (70,-11)
      {바깥 고리의 출력이\\곧 안쪽 고리의 입력임};
\node[note, anchor=south, align=center, font=\scriptsize] at (240,-6)
      {VRX 쌍둥이는 이 칸이\\Gazebo 로 바뀜};

\end{tikzpicture}
\end{document}
